\subsection{Reduction to an auxiliary permanent}\label{costes:permanente}

Let $\varepsilon_{\rm out}\ge\max_f\ell_f$ be the common outer bound, with $\varepsilon_{\rm out}<1$.
Let $P=(p_{ij})$ be the matrix of the interior stationary Bethe point, with zeros off the support, and put
$W_{ij}=p_{ij}(1-p_{ij})$. In what follows $W$ denotes this auxiliary matrix. Write
\begin{equation}\label{coste:rook-def}
 \begin{gathered}
 S_{R,d}=A_R^T\operatorname{diag}(\tau_d(p_e))A_R,\qquad S_2=S_{R,2},\\
 C_2=\sum_{ij}\bigl[\log(1-2p_{ij})-2\log(1-p_{ij})\bigr],\qquad
 \Gamma_2=\frac{e^{-C_2/2}\sqrt{4\pi k}}{\sqrt{\pdet S_2}} .
 \end{gathered}
\end{equation}
A permutation on the support uses one square in each row and each column. We call it legal
if it also uses each diagonal at most once.

\begin{lema}[Reduction to a permanent]\label{coste-factor}
Let $\mu_W(\sigma)=\prod_{e\in\sigma}W_e/\per W$ be the probability that $W$ assigns to a
permutation on the support, and let
\[
 F(\sigma)=
 \frac{\prod_f h_f}{\prod_e(1-p_e)^2}\,
 \mathbf1_{\{\sigma\ {\rm legal}\}}
 \prod_{e\in\sigma}\frac{(1-p_e)^2}{h_{D(e)}h_{A(e)}} .
\]
Then
\[
 Q(B)e^{-H_B(p)}
   =\frac{\per W}{\prod_e(1-p_e)}\,\E_{\mu_W}F.
\]
For $Q(B)>0$, suppose the doubly stochastic matrix $P$ satisfies the bounds
in \eqref{terminal-hipotesis-matriz} with $K=k$, namely $p_{ij}\le M_0/k$ and
$\|P-\mathbf1\mathbf1^T/k\|_{\rm op}\le\rho_*<1$. Then uniformly for fixed $M_0,\rho_*$,
\[
 g(B)=\log\E_{\mu_W}F-\log(\Gamma_4/\Gamma_2)+o(1)
 \qquad (k\longrightarrow\infty).
\]
\end{lema}
\begin{proof}
The rows and columns of $P$ sum to one. The mass of any set of rows is contained in its
neighboring columns, each of capacity one. Hall's theorem gives a permutation on the support,
so $\per W>0$. On a legal permutation, multiplying $W_e$ by the fraction in the last product gives
\[
 \frac{p_e(1-p_e)^3}{h_{D(e)}h_{A(e)}}=a_i b_j
\]
by Lemma~\ref{coste-vacancias}. The product over the permutation is therefore $\prod_i a_i\prod_j b_j$,
independently of the placement. After summing over the $Q(B)$ legal permutations, the right-hand
side of the first equality is $Q(B)\mathcal K/\prod_e(1-p_e)^3$, which equals $Q(B)e^{-H_B(p)}$
by the same lemma.

For the asymptotic formula apply the matrix estimate of Lemma~\ref{terminal-normalizador} to this
$P$, with $K=k$, $W_{ij}=p_{ij}(1-p_{ij})$, and the rook factor \eqref{coste:rook-def}. It gives
$\log(\per W/\prod_e(1-p_e))=\log\Gamma_2+o(1)$. Its proof uses the scaling and permanent
estimates of Section~\ref{permanente}. The multipliers in the weighted rook stationarity equation
$p_e(1-p_e)=W_e$ can all be taken equal to one. Taking logarithms in the first equality and using
$L(B)=H_B(p)+\log\Gamma_4$ proves the second. We use the law $\mu_W$ only for counting on $B$.
\end{proof}

\subsection{Terms in the ratio of normalizers}

\begin{lema}[Inverses of the row and column matrices]\label{coste-inversa}
Under the hypotheses of Lemma~\ref{coste-estatico}, write $S_0=A_R^T\operatorname{diag}(p)A_R$
and let $\iota_e$ be the indicator vector of the row and column of $e$.
For $0\le s\le1$ put $S(s)=(1-s)S_{R,2}+sS_{R,4}$. Uniformly in $s$ and in the physical indices
$u,v$,
\[
 (S(s)^+)_{uv}=\mathbf1_{\{u=v\}}+O(k^{-1}).
\]
Consequently $\iota_e^TS(s)^+\iota_{e'}$ equals the number of row and column vertices shared
by $e,e'$, up to an error of $O(k^{-1})$.
\end{lema}
\begin{proof}
Since the row and column masses equal one,
\[
 S_0=\begin{pmatrix}I&P\\P^T&I\end{pmatrix}.
\]
Let $\Pi_+,\Pi_-$ denote the orthogonal projections onto the two lines spanned by
$(\mathbf1_R,\mathbf1_C)$ and $(\mathbf1_R,-\mathbf1_C)$, respectively. Their entries have absolute
value $1/(2k)$. The corresponding eigenvalues are two and zero. Let $N_0$ be the symmetric matrix
with zero diagonal blocks and off-diagonal blocks $P-\mathbf1\mathbf1^T/k$ and its transpose.
It annihilates both lines of constant vectors, and on their orthogonal complement $S_0=I+N_0$.
The rows and columns of $P$ have Euclidean norm $O(k^{-1/2})$ by $\sum_jp_{ij}^2\le C/k$ and its
column analogue. The same holds for the rows of $N_0$. Its entries are $O(1/k)$, and
$\|(I+N_0)^{-1}\|_{\rm op}\le1/\min\{1,\gamma\}$. The exact identity
\[
 (I+N_0)^{-1}= I-N_0+N_0(I+N_0)^{-1}N_0
\]
and Cauchy--Schwarz applied to two rows of $N_0$ show that the entries of the inverse differ from
those of $I$ by $O(k^{-1})$. To recover $S_0^+$, subtract $\Pi_-+\tfrac12\Pi_+$. This preserves the
same estimate.

For $d=2,4$ we have $0\le\tau_d(p_e)-p_e=O(p_e^2)$. In each row of $S_{R,d}-S_0$ the sum of
absolute values is twice the sum of these differences along one physical line and is $O(k^{-1})$.
By symmetry the same bound holds for the operator norm, also for $S(s)-S_0$. All these matrices
have the same kernel and dominate $S_0$ on its orthogonal complement. On this complement, the
resolvent identity gives
$\|S(s)^+-S_0^+\|_{\rm op}\le\gamma^{-2}\|S(s)-S_0\|_{\rm op}=O(k^{-1})$.
Contracting with $\iota_e,\iota_{e'}$ involves only four entries. This gives the remaining claim.
\end{proof}
